\section{Full Proofs for Section~\ref{sec:pof}}

\begin{proof}[of Theorem~\ref{thm:pof}]
\emph{Endpoint.} If $\theta=0$, repair is unavailable; dropping every OPT
batch to the singleton floor allocation loses at most the factor $g$, and
the lower-bound construction below attains $g$. Hence assume $\theta>0$
for the repair calculation.

\emph{Upper bound.} Fix an unconstrained optimum and process its winning
bids independently; bids cover disjoint order sets, so the repaired
welfare contributions add. Take a winning bid with floor mass $C_S=\sum_{j\in S}c_j>0$
(a bid with $C_S=0$ has all floors zero and is already fair), value
$V=\sum_{j\in S}y_j$, where $y_j:=y_{i,S,j}$, gain
$\gamma=V/C_S$, and deficit $D=\sum_{j\in S}(c_j-y_j)^+$. By local
optimality of $\OPT$, we may assume $V\ge C_S$ (otherwise replacing this
bid by the singleton floor allocation on $S$ improves welfare), so
$\gamma\ge1$; Definition~\ref{def:g} gives $\gamma\le g$. Two fair uses:

\emph{Repair.} Move value from over-served legs to starved ones. The
total surplus available is $s=V-(C_S-D)=V-C_S+D$, and clearing the
deficit costs $D/\theta$ units of surplus (haircut $\lambda=1-\theta$).
This is feasible iff $V-C_S\ge D\frac{1-\theta}{\theta}$, and then the
delivery saturating the conversion budget of Definition~\ref{def:theta}
has welfare $f=V-D\frac{1-\theta}{\theta}$. Since every raw output
value is nonnegative, $(c_j-y_j)^+\le c_j$ and therefore $D\le C_S$.
Together with $V-C_S=(\gamma-1)C_S$, feasibility holds whenever
$\gamma-1\ge\frac{1-\theta}{\theta}$, i.e. $\gamma\ge1/\theta$; then
\[
  \frac{f}{V}
  =1-\frac{1-\theta}{\theta\gamma}\frac{D}{C_S}
  \ge 1-\frac{1-\theta}{\theta\gamma}
  \ge \theta
  = \frac{1}{1/\theta}.
\]
Thus repair gives $f\ge V/(1/\theta)$ whenever $\gamma\ge1/\theta$, and
hence $f\ge V/\!\min(g,1/\theta)$ in the repair regime.

\emph{Drop.} Settle every covered order at its floor: $f=C_S=V/\gamma$.
For $\gamma\le\min(g,1/\theta)$ this is $\ge V/\!\min(g,1/\theta)$.

Either way each bid retains a $1/\min(g,1/\theta)$ fraction of its value;
summing, $\FAIR\ge\OPT/\min(g,1/\theta)$.

\emph{Lower bound.} Let $\gamma^\ast=\min(g,1/\theta)$. Two orders on
distinct pairs, $c_1=0$, $c_2=c$; one batched bid with
$y_2=0$, $y_1=\gamma^\ast c$, so $V=\gamma^\ast c$, $C_S=c$, and the
batching solver's standalone values are $0,c$, so
Definition~\ref{def:g} holds as $\gamma^\ast\le g$. The unconstrained
optimum runs this batch, delivering $\gamma^\ast c$ (all on order~1), so
$\OPT=\gamma^\ast c$. Fair repair of order $2$ to its floor $c$ requires
withdrawing $c/\theta$ from $y_1$; if $\gamma^\ast=1/\theta$ this exhausts
$y_1$ exactly (delivered welfare $c$), and if
$\gamma^\ast=g<1/\theta$ the batch is irreparable (welfare $c$ from the
individual allocation). In both cases $\FAIR=c$ and
$\OPT=\gamma^\ast c$, so $\OPT/\FAIR=\gamma^\ast$. Independence of the
number of orders and of prices is immediate from the construction.\hfill$\square$
\end{proof}

\begin{proof}[of Proposition~\ref{prop:dirrel}]
The two-order family of Theorem~\ref{thm:pof} already has $d=2$ and
attains $\min(g,1/\theta)$, and the upper bound never uses $d$; for
$d=1$ only singletons exist, so $\OPT=\FAIR=C$.\hfill$\square$
\end{proof}

\begin{proof}[of Theorem~\ref{thm:coverage}]
\emph{Endpoint.} If $\theta=0$ and $\alpha<1$, repair is unavailable and
the drop argument gives ratio $g$; if $\alpha=1$, every admitted bid
already clears all floors and the ratio is $1$. Hence assume $\theta>0$
below.

\emph{Upper bound.} Write
$\Gamma=\Gamma(\theta,\alpha):=\alpha+(1-\alpha)/\theta$, with the
endpoint convention from the theorem. As above, coverage caps the deficit by
$D\le(1-\alpha)C_S$. Repair is feasible when
$\gamma-1\ge(1-\alpha)\frac{1-\theta}{\theta}=\Gamma-1$, and then
$f=V-\frac{D(1-\theta)}{\theta}\ge V\big(1-\frac{(1-\alpha)(1-\theta)}
{\theta\gamma}\big)\ge V/\Gamma$ for $\gamma\ge\Gamma$ (monotone in
$\gamma$, equality at $\gamma=\Gamma$); dropping gives $V/\gamma\ge
V/\Gamma$ for $\gamma\le\Gamma$ and $V/g$ for $\gamma\le g$. Hence
$\FAIR\ge\OPT/\min(g,\Gamma)$. \emph{Lower bound.} Two orders, $c_1=0$,
$c_2=c$; batched bid with $y_2=\alpha c$ (coverage saturated) and
$y_1=(\gamma^\ast-\alpha)c$, $\gamma^\ast=\min(g,\Gamma)$; the batching
solver's standalone values are $0,c$, so Definition~\ref{def:g} holds as
$\gamma^\ast\le g$. Repairing
order $2$ needs $(1-\alpha)c/\theta$ from $y_1$; if $\gamma^\ast=\Gamma$
then $y_1=(\Gamma-\alpha)c=\frac{(1-\alpha)c}{\theta}$ exactly, and if
$\gamma^\ast=g<\Gamma$ the batch is irreparable. Either way $\FAIR=c$,
$\OPT=\gamma^\ast c$.\hfill$\square$
\end{proof}

\section{Full Proofs for Section~\ref{sec:walras}}

\begin{proof}[of Lemma~\ref{lem:domination}]
\[
\begin{aligned}
\sum_{j\in S}(y_{i,S,j}-w_j)
&=V_i(S)-\sum_{j\in S}w_j\\
&\le \sum_{j\in S}V_i(\{j\})-\sum_{j\in S}w_j\\
&\le \sum_{j\in S}(V_i(\{j\})-w_j)^+ .
\end{aligned}
\]
The middle step uses $g=1$. Each
positive-profit singleton has $V_i(\{j\})>w_j$, hence is executable; a
non-executable batch is not a deviation. The obligation $w_j$ is the same
whether $j$ is served in the batch or individually, which is the only
property of UDP used.\hfill$\square$
\end{proof}

\begin{proof}[of Theorem~\ref{prop:walras}]
\emph{$(a)$, $g=1$.} Set $w_j=c_j$ and assign each order $j$ to a solver
attaining $V_i(\{j\})=c_j$. The winner earns zero margin per order; any
losing singleton has $V_i(\{j\})-c_j\le0$; and
Lemma~\ref{lem:domination} kills batch deviations (no positive-profit
singletons). With $g=1$, any allocation produces at most
$\sum_{(i,S)}V_i(S)\le\sum_j c_j=C$, attained here, so produced value is
$C=\OPT$, all delivered. The point $w_j=c_j$ is the best-bid floor, so
the equilibrium is fair (the trader-optimal core point of the underlying
assignment market~\cite{shapley1971}, to which
Lemma~\ref{lem:domination} reduces the economy).

\emph{$(b)$, $g>1$.} We argue natively; non-existence with transfers does
not carry over, as feasibility only removes deviations. Fix
$\delta\in(0,2(g-1)]$; orders $1,2,3$ on distinct pairs, so the floor
vector $c_j=1$ is rate-realizable. A rival has
standalone value $1$ on each ($c_j=1$), and three solvers each holding
one batched bid on a pair $\{j,j'\}$ with vector
$(1+\tfrac\delta2,1+\tfrac\delta2)$ (value $2+\delta\le2g$) and
standalone values $1$. The efficient allocation runs one batch, say
$\{1,2\}$, plus order $3$ individually (produced $3+\delta$). Suppose
rates $w$ support it. Order $3$ is served by a standalone solver of value
$1$, so $w_3\le1$; orders $1,2$ are served by the batch, so
$w_1,w_2\le1+\tfrac\delta2$. The batch $\{2,3\}$ is then executable, and
equilibrium requires it unprofitable against its solver's individual
alternatives: $(2+\delta)-w_2-w_3\le(1-w_2)^++(1-w_3)^+$. If $w_2\le1$
the right side is $(1-w_2)+(1-w_3)$, giving $\delta\le0$; if $w_2>1$ it
is $1-w_3$, giving $w_2\ge1+\delta>1+\tfrac\delta2$, contradicting
feasibility. No supporting rates exist.\hfill$\square$
\end{proof}

\begin{proof}[of Proposition~\ref{prop:core}]
(i) If $w_j<c_j$, the standalone bid attaining $c_j$ blocks $j$ at factor
$1$; if $w_j\ge c_j$ for all $j$, no singleton strictly improves. (ii)
Blocking $w_j=c_j$ at factor $g$ needs delivery $>g\,c_j$ to every
$j\in S$, totaling $>g\,C_S$, while any feasible delivery from a bid
totals $\le V_i(S)\le g\,C_S$ (haircuts only lose value; if all floors in
$S$ are zero then $V_i(S)=0$ and strict blocking is impossible).
Tightness: a batch with vector $(g'c,g'c)$ over floors $(c,c)$,
$g'\in(\kappa,g]$, blocks at any $\kappa<g$. (iii) Three-cycle floors
$(1,1,1)$, pair vectors $(1+\tfrac\delta2,1+\tfrac\delta2)$: summing the
three transferable core constraints $x_j+x_{j'}\ge2+\delta$ gives
$\sum x_j\ge3+\tfrac{3\delta}2>3+\delta=\OPT$, so the transferable core is
empty; the allocation (batch $\{1,2\}$, order $3$ individually) is
$\theta=0$-unblocked since every other batch caps a covered order at a
value already attained.\hfill$\square$
\end{proof}

\paragraph{Complexity.}
NP-hardness is a reduction from 3-dimensional matching: orders are ground
elements, each element has an individual bid of value $1$, and each triple
is represented by a solver whose singleton values on the three covered
elements are all $1$ and whose triple bid produces $1+\delta/3$ on each of
the three orders (total value $3+\delta\le3g$ for $\delta\le3(g-1)$). Thus
the triple satisfies the complementarity bound and passes the fairness
filter.
If the maximum number of disjoint triples is $m$, the optimum welfare is
$n+\delta m$, where $n$ is the number of ground elements; hence welfare
$n+\delta n/3$ is attainable iff a perfect matching exists. For
fixed-parameter tractability, treat each directed pair as an atomic item
under UDP: all orders on the same directed pair settle together, so a
feasible solution either assigns the entire directed-pair block to one
bid or leaves it unserved. After fair filtering, every surviving bid is a
weighted subset of the $P$ atoms, and winner determination is weighted set
packing on a universe of size $P$. The dynamic program
$DP[M]=\max\{DP[M\setminus T_b]+w_b:b\text{ uses }T_b\subseteq M\}$
over masks $M\subseteq[P]$ scans the bid list at each state, giving
$2^{O(P)}\cdot\mathrm{poly}(|\mathcal B|)$.

\section{Full Proofs for Section~\ref{sec:incentives}}

\begin{proof}[of Theorem~\ref{thm:incompat}]
Let there be one order. Solver $i$ has private deliverable value $v_i$
and may report any binding value $r_i\le v_i$; over-reporting is excluded
by settlement, while under-reporting is feasible. By non-wastefulness the
order is allocated whenever some report is positive, and by
value-respecting the winner has a maximal report. Feasibility implies the
winner can deliver at most its true value. Best-bid fairness requires the
delivery to be at least the endogenous benchmark $\max_i r_i$.

Let the highest and second-highest true values be
$v^{(1)}>v^{(2)}$. If the top solver reports truthfully, it sets the
benchmark to $v^{(1)}$, must deliver $v^{(1)}$, and earns zero surplus.
If instead it reports $r=v^{(2)}+\varepsilon$ with
$0<\varepsilon<v^{(1)}-v^{(2)}$, it still has the maximal report, the
benchmark falls to $v^{(2)}+\varepsilon$, and it can earn
$v^{(1)}-v^{(2)}-\varepsilon>0$. Truthful reporting is therefore not a
dominant strategy for any feasible non-wasteful value-respecting
mechanism enforcing the best-bid benchmark. The same argument shows that
any DSIC mechanism in this class can enforce at most the runner-up
benchmark.\hfill$\square$
\end{proof}

\begin{proof}[of Theorem~\ref{thm:posp}]
\emph{Endpoint.} If $g=1$, the floor and the batch value coincide, so
settling at the floor is optimal and strategyproof. Hence assume $g>1$
for the lower-bound expression.

\emph{Lower bound.} Restrict to the screening family $R_g$: a single
strategic solver holds one batch of private value $V\in[1,g]$ (floor mass
normalized to $1$, so $\FAIR(V)=V$), and the mechanism posts an outcome
as a function of the reported $V$. Write $p(V)$ for the probability the
batch is executed and $x(V)$ for the expected batch delivery (zero when
not executed). The fallback floor delivers welfare $1$, so
$W(V)=x(V)+1-p(V)$, and the solver's interim utility is
$u(V)=p(V)V-x(V)$. Because bids are binding, a type $V$ can only
\emph{under}report (claim $V'\le V$), so incentive compatibility implies
\[
u(V)\ge u(V')+(V-V')p(V') \qquad (V'\le V).
\]
Thus $u$ is nondecreasing and $u'(V)\ge p(V)$ a.e. Hence
$u(V)\ge\int_1^V p(t)\,dt$ and
\[
x(V)=p(V)V-u(V)\le
\psi(V):=p(V)V-\int_1^V p(t)\,dt .
\]
Integrating $\psi(V)/V^2$ over $[1,g]$ and applying Fubini gives
\[
\int_1^g\frac{\psi(V)}{V^2}\,dV
=\int_1^g\frac{p(V)}{V}\,dV
-\int_1^g p(t)\Big(\frac1t-\frac1g\Big)\,dt
=\frac1g\int_1^g p(t)\,dt
\le\frac{g-1}{g}.
\]
If the mechanism has competitive ratio $r$, i.e. $W(V)\ge V/r$
pointwise, then $x(V)\ge V/r-(1-p(V))\ge V/r-1$, and
\[
\int_1^g \frac{V/r-1}{V^2}\,dV
=\frac{\ln g}{r}-\Big(1-\frac1g\Big).
\]
Combining the two inequalities,
$(\ln g)/r\le2(1-1/g)$, so every randomized DSIC mechanism has
\[
r=\sup_{V\in[1,g]}\frac{V}{W(V)}
\ge\frac{g\ln g}{2(g-1)}
\ge\frac{\ln g}{2}.
\]

\emph{Posted-price upper bound.} Offer the batch a take-it-or-leave-it
delivered welfare on the geometric ladder $\{e^k:0\le k\le K\}$,
$K=\lceil\ln g\rceil$. Choose $k$ uniformly at random, execute the batch
if it accepts the posted threshold, and otherwise settle at the floors.
For any realized $V\in[e^k,e^{k+1})$, the draw $k$ (probability
$1/(K+1)$) executes the batch and delivers at least $e^k>V/e$, while all
draws deliver at least the floor mass. Hence
$\E[\text{welfare}]\ge\max\big(V/(e(K+1)),1\big)\ge\FAIR/(e(\lceil\ln
g\rceil+1))$.

\emph{Deterministic optimum.} Let the mechanism be deterministic. First,
its execution rule is a threshold: if $p(r)=1$ and $p(V)=0$ for some
$r<V$, type $V$ has truthful utility $0$, while under-reporting $r$ yields
utility $V-x(r)\ge V-r>0$ (feasibility at report $r$ gives $x(r)\le r$),
contradicting downward IC. Thus there is a threshold $\tau$ (up to the
boundary point) such that types below $\tau$ are not executed and types
above $\tau$ are. Second, the accepted delivery is bounded by $\tau$: for
$\tau\le r<V$, downward IC gives $V-x(V)\ge V-x(r)$, hence
$x(V)\le x(r)$; letting $r\downarrow\tau$ and using feasibility
$x(r)\le r$ gives $x(V)\le\tau$ for all executed types, in particular
$x(g)\le\tau$. Types just below $\tau$ have welfare $1$, while type $g$
has welfare at most $\tau$, so every deterministic DSIC rule has ratio at
least $\max\{\tau,g/\tau\}\ge\sqrt g$. The posted delivery $\sqrt g$
attains this bound. Since $R_g$ embeds as a degenerate single-minded
subfamily of the full model, all three bounds transfer.\hfill$\square$
\end{proof}

\begin{proof}[of Theorem~\ref{thm:lift}]
\emph{(i) Eligibility-posted packing.} In the public-floor regime, draw
$k\sim\mathrm{Unif}\{0,\dots,K\}$, $K=\lceil\ln g\rceil$, post the weight
$C_S e^{k}$ for each option $S$, and let a solver accept its
option iff its binding value clears this weight; among accepted options
take a maximum-weight disjoint packing using the posted weights. Accepted
options deliver their posted weights componentwise by the filter-passing
scaling condition from the mechanism model; all orders not covered by the
selected posted options settle at their public floors. The
posted weights and packing rule are report-independent, so each solver
faces a take-it-or-leave-it offer and the mechanism is universally
truthful. Let
$X$ be the total value of the batch options used by the offline fair
optimum. The geometric draw captures expected posted delivery at least
$X/(e(K+1))$. Independently, every realization settles all orders at least
at their public floors, so $W\ge C$. Since $\FAIR\le X+C$, the
max-inequality with $a=e(K+1)$ gives
$\E[W]\ge\FAIR/(e(\lceil\ln g\rceil+1)+1)$.

\emph{(ii) Coupled menu.} Draw $k$ uniformly from $\{0,\ldots,K\}$ and
set $t=e^k$. Post every option at price $C_S\,t$. Let the solver demand a
profit-maximizing disjoint collection $\mathcal C^\ast(t)$. This is
universally truthful because the posted menu is report-independent. It is
feasible because any demanded $S$ has $V_S\ge C_St\ge C_S$, and the
filter-passing scaling condition from the mechanism model implements total
delivery $C_St$ while clearing floors componentwise. Let
\[
p(t)=
\max_{\mathcal C\text{ disjoint}}
\sum_{S\in\mathcal C}(V_S-C_St).
\]
As a maximum of finitely many affine decreasing functions, $p$ is convex,
piecewise linear, nonincreasing, with $p(t)=0$ for $t\ge g$ (every
option has $V_S\le gC_S$). At any differentiability point
$-p'(t)=\sum_{S\in\mathcal C^\ast(t)}C_S$, so the delivered mass is
$D(t)=t\sum_{S\in\mathcal C^\ast(t)}C_S=t\,(-p'(t))$; at a kink the
right-derivative selects the continuing (smallest-floor) active
collection, and $D(t)\ge t\,(-p'_+(t))$. On $[t_k,t_{k+1}]=[e^k,e^{k+1}]$,
convexity gives $-p'$ nonincreasing, so
\[
\int_{t_k}^{t_{k+1}}(-p'(t))\,dt\le (-p'_+(t_k))(t_{k+1}-t_k)
=(-p'_+(t_k))(e-1)e^k\le (e-1)D(t_k).
\]
Summing over $k=0,\dots,K-1$ and using $p(e^K)=0$,
\[
p(1)=p(1)-p(e^K)=\int_1^{e^K}(-p'(t))\,dt\le(e-1)\sum_{k=0}^{K}D(e^k),
\]
so the uniform draw gives $\E[D]\ge p(1)/((e-1)(K+1))$. Since the fair
optimum's batches form a feasible disjoint collection at $t=1$,
$p(1)\ge\sum_{S\in B^\ast}(V_S-C_S)=\FAIR-C$, and combining with the
pointwise $W\ge C$ via the same $\max$-inequality yields
$\E[W]\ge\FAIR/((e-1)(\lceil\ln g\rceil+1)+1)$.

\emph{Independent geometric pricing fails.}
The separation uses the following construction. There are options
$S_t=\{0,t\}$ for $t=1,\dots,m$, all sharing the common order $0$, with
$c_0=c_t=\tfrac12$ and $C_{S_t}=1$. The solver has value $V_t=g$ for every
option. Let $P_t$ be the posted price of option $S_t$ and
$P_{\min}:=\min_t P_t$. Independent geometric prices make the solver
choose only the cheapest acceptable option: $D=P_{\min}$ if
$P_{\min}\le g$, and $D=0$
otherwise. Hence
\[
D\le1+g\,\mathbf 1\{P_{\min}\ge e\}.
\]
Moreover,
\[
\Pr[P_{\min}\ge e]=(K/(K+1))^m\le e^{-m/(K+1)} .
\]
Choosing $m=\lceil(K+1)\ln g\rceil$ gives $\E[D]\le2$. Thus
$\E[W]\le2+C=O(\log^2 g)$, while the fair optimum obtains
$\FAIR=g+(m-1)/2=\Omega(g)$ by taking one high-value option and settling
the remaining leaf orders at floors. Hence the independent geometric
analogue has ratio $\Omega(g/\log^2 g)$.\hfill$\square$
\end{proof}

\begin{proof}[of Theorem~\ref{thm:multi}]
\emph{Upper bound.} In random-priority coupled menus, condition on a
solver being served first (probability $1/N$): it faces the full coupled
menu on all orders, and the profit-curve telescope of
Theorem~\ref{thm:lift}(ii) applies verbatim, giving expected delivery at
least $p_i(1)/((e-1)(K+1))$ from that solver, where $p_i$ is computed on
the full order set. A solver served later faces a subset of orders and
contributes nonnegatively. Since each solver's optimal batches form a
disjoint collection on the full set, $\sum_i p_i(1)\ge\sum_{b\in
B^\ast}(V_b-C_b)=\FAIR-C$. Averaging over the uniform priority and
combining with $W\ge C$ gives
$\E[W]\ge\FAIR/(N(e-1)(\lceil\ln g\rceil+1)+1)$. Truthfulness holds
because a solver's menu depends only on which orders predecessors have
consumed, never on its own report.

\emph{Lower bound (symmetrized swarm).} For $g=1$ the lower bound is the
trivial constant bound, so assume $g>1$. Fix orders $o_1,\dots,o_N$ with
floors $\varepsilon=L/(N\rho)$ and auxiliary orders $u_1,\dots,u_L$ with
floor $1$; let $\mathcal U$ be the set of all orders, so
$C_{\mathcal U}\le2L$. Each identity
has the \emph{same} public option family
$\{\{o_1\},\ldots,\{o_N\},\mathcal U\}$; identities
differ only in private values, so any
report-independent service order is statistically independent of which
identity is the ``batcher.'' Choose $\rho\in(\max\{1,g/2\},g]$ avoiding
the (countably many) atoms of the $N$ multiplier distributions. The
batcher, placed by the adversary at an identity chosen after seeing the
mechanism, has $V(\mathcal U)=\rho\,C_{\mathcal U}$ and ratio-$1$
singletons. Every other identity (blocker $i$) has a ratio-$\rho$
singleton on $\{o_i\}$ and ratio-$1$ on all other options, including
$\mathcal U$. At multiplier $t$ a blocker is silent iff $t>\rho$;
write $q_i=\Pr[t_i>\rho]$. The batcher delivers
$\Theta(\rho C_{\mathcal U})$ only if
it is reached before any blocker has consumed an order of $\mathcal U$, i.e. only
if all predecessors are silent. Blockers together deliver at most
$N\varepsilon\rho=L$, and fallback floors contribute at most
$C_{\mathcal U}\le2L$,
so pointwise $W\le\mathrm{del}_{\mathrm{batcher}}+3L$. For any
realization of the order,
\[
\sum_{i}(1-q_i)\!\!\prod_{s\prec i}\!q_s \;=\; 1-\prod_s q_s \;\le\; 1,
\]
the $i$-th term being the probability that $i$ is the first non-silent
identity. Taking expectation over the (independent) order and the
multipliers,
\[
  \sum_i (1-q_i)\Pr[\text{predecessors of }i\text{ silent}] \le 1,
\]
so some identity $i^\ast$ has
$(1-q_{i^\ast})\Pr[\text{predecessors silent}]\le 1/N$. Placing the
batcher at $i^\ast$ caps its expected batch delivery at $2\rho L/N$,
while the fair optimum delivers at least $\rho L$ via the batch. Hence
\[
\frac{\FAIR}{\E[W]}\ge\frac{\rho}{3+2\rho/N}
=\Omega(\min(N,\rho))=\Omega(\min(N,g)),
\]
using $\rho>g/2$. The blockers' geometry is public and identical across
placements, so the bound holds even for geometry-aware orders.\hfill$\square$
\end{proof}

\section{Full Proofs for Section~\ref{sec:online}}

\begin{proof}[of Proposition~\ref{prop:online}]
Settle-at-deadline is fair and achieves $C$ on every instance, while
$\FAIR\le\OPT\le gC$, so the ratio is at most $g$; a staggered-deadline
instance forcing individual-only settlement shows it is at least $g$.\hfill$\square$
\end{proof}

\begin{proof}[of Theorem~\ref{thm:online}]
(i) If an almost-surely fair mechanism settles an exposed batch (a
covered order's deadline in the future) with positive probability on some
prefix, fix that prefix and append a higher standalone bid on a
later-deadline covered order: on the appended instance it has, with
positive probability, settled below a realized floor---not a.s.\ fair.
Hence exposed batches are never settled, and on the staggered hard family
the mechanism is individual-only a.s., with welfare $\le C$ against
$\FAIR=gC$.
(ii) Stopped welfare $q=V+C-C_S$ satisfies $q\in[C,gC]$ ($V\ge C_S$ and
$V\le gC_S$, $C_S\le C$). \emph{Deterministic} $\Theta(\sqrt g)$: accept
the first $q\ge\sqrt g\,C$; if $q^\ast\ge\sqrt g\,C$ the ratio is
$q^\ast/q\le\sqrt g$, else $W=C$ and ratio $<\sqrt g$; a two-step
adversary matches. \emph{Randomized} $O(1+\log g)$: with
$k\sim\mathrm{Unif}\{0,\dots,K\}$, accept the first $q\ge e^kC$. With
probability $1/(K+1)$, $e^kC\le q^\ast<e^{k+1}C$, and then the mechanism
accepts some offer with $q\ge e^kC>q^\ast/e$, so
$\E[W]\ge q^\ast/(e(K+1))$. The matching $\Omega(1+\log g)$ bound is
constant at $g=1$; for $g>1$ its logarithmic term is a Yao bound:
a persistent order with fleeting geometric batch offers $x_i=e^ic$,
$i\le\lfloor\ln g\rfloor$, shown for the first $M$ with
$\Pr[M\ge i]=e^{-i}$, caps every deterministic threshold at $2c$ while
$\E[\FAIR]=\Omega(c\log g)$. Greedy: a ratio-$g$ blocker arriving before a
ratio-$g$ batch is accepted (killing the batch) for any threshold
$T\le g$ and rejected with the batch for $T>g$, ratio $\Omega(g)$.\hfill$\square$
\end{proof}

\begin{proof}[of Theorem~\ref{thm:poa}]
In the per-order-decomposable individual game, the maximal standing bid
on $j$ both wins and delivers itself, so $c_j(b)=p_j(b)$. Let
$v_j^{\max}:=\max_i v_{i,j}$, choose
$i^\ast(j)\in\arg\max_i v_{i,j}$, and set
$\OPT_{\mathrm{ind}}:=\sum_j v_j^{\max}$. Let $\beta_j$ denote the
price/penalty term faced by the deviation on order $j$; below
$\beta_j\le c_j(b)=p_j(b)$. Let
$i^\ast(j)$ deviate to $W_j$ with density $1/(v^{\max}_j-w)$ on
$[0,(1-1/e)v^{\max}_j]$ (Syrgkanis--Tardos). Then the expected utility
from $j$ is $((1-\tfrac1e)v^{\max}_j-\beta_j)^+\ge(1-\tfrac1e)
v^{\max}_j-p_j(b)$ with $\beta_j\le c_j(b)=p_j(b)$, the penalty landing
under $b$. Summing and using
$\sum_iu_i(b)=\PW(b)-\sum_jp_j(b)$ in the CCE inequality cancels the
price terms, giving $\E[\PW]\ge(1-1/e)\OPT_{\mathrm{ind}}$; against the
unrestricted optimum, $\OPT\le g\,\OPT_{\mathrm{ind}}$, so the factor is
at most $(e/(e-1))g$.
The full-game obstruction is a latent batch present in $b_{-i}$
that executes under the deviation but not under $b$, breaking every
charging scheme.\hfill$\square$
\end{proof}

\section{Benchmark Manipulation}\label{app:manipulation}

On one directed pair, let an attacker of volume $v$ depress a quantile
benchmark by injecting mass $a$ at convex cost
$K(a)=\kappa_1a+\tfrac{\kappa_2}2a^2$, with local sensitivity $\beta$.
For a lower-tail quantile of $m$ honest quotes with density $f$ at
$b^\ast$, first-order sensitivity is $\beta=1/(mf(b^\ast))$.
The attacker chooses $a_t\in[0,\bar a]$ each round, with $\bar a>0$ and
$\kappa_2>0$. A public rate limit $r\in(0,1)$ bounds the downward bias
after a restart from $b_0=b^\ast$ by $b^\ast(1-(1-r)^t)$ in round $t$.
Steady-bias statements assume $\mu<1$.

\begin{theorem}[The contamination game]\label{thm:contam}
A marginal injection cost $\kappa_1\ge v\beta$ deters contamination.
Otherwise the optimal steady attack is
$a^\ast=\min\{\bar a,(v\beta-\kappa_1)/\kappa_2\}$, inducing relative
bias $\mu=\beta a^\ast/b^\ast$; bursts do not beat the constant attack,
and rate limits tax restarts through a finite ramp without changing the
steady state. Every public-floor guarantee is then read with
$G\mapsto G/(1-\mu)$.
\end{theorem}

\begin{proof}[of Theorem~\ref{thm:contam}]
The steady-state benchmark under a constant injection $a$ is
$b^\ast-\beta a$ (the rate limit does not bind once the bias is
constant), so the per-round profit is
$\pi(a)=v\beta a-\kappa_1 a-\tfrac{\kappa_2}{2}a^2$. With $\kappa_2>0$,
$\pi$ is strictly concave with $\pi'(0)=v\beta-\kappa_1$; hence if
$v\beta\le\kappa_1$ the unique maximizer is $a^\ast=0$ (no attack), and
otherwise $a^\ast=\min\{\bar a,(v\beta-\kappa_1)/\kappa_2\}$, with
relative bias $\mu=\beta a^\ast/b^\ast$.

\emph{Bursts do not help.} The published benchmark satisfies
$b_t\ge b^\ast-\beta a_t$ pointwise, so for every horizon $T$ and every
(possibly time-varying) policy $(a_t)$,
\[
\sum_{t=1}^T\big[v(b^\ast-b_t)-K(a_t)\big]
\le \sum_{t=1}^T\big[v\beta a_t-K(a_t)\big]
=\sum_{t=1}^T\pi(a_t)\le T\,\pi(a^\ast),
\]
and the constant policy $a^\ast$ attains $T\pi(a^\ast)-O(1)$ after the
ramp; hence no episodic schedule beats the constant attack.

\emph{Ramp tax (exact).} Starting from $b_0=b^\ast$, the rate limit caps
the attainable bias at round $t$ by $\Delta_t:=\min\{\Delta^\ast,\,
b^\ast(1-(1-r)^t)\}$ with $\Delta^\ast=\beta a^\ast$; reaching
$\Delta^\ast$ takes $T_\mu=\big\lceil\ln\tfrac{1}{1-\mu}\big/
\ln\tfrac{1}{1-r}\big\rceil$ rounds. An attacker who hugs the cap injects
$a_t=\Delta_t/\beta$, so the exact net rent forgone per restart is
\[
L_{\mathrm{net}}=\sum_{t=1}^{T_\mu}\Big[\pi(a^\ast)-\pi(\Delta_t/\beta)\Big]
\;\le\; \pi(a^\ast)\,T_\mu .
\]
When the cap is slack, so $a^\ast<\bar a$, and in the small-$(r,\mu)$
regime, $1-(1-r)^t\approx rt$ and the summand is
$\tfrac{\kappa_2}{2}(a^\ast)^2(1-rt/\mu)^2$, giving
$L_{\mathrm{net}}=\tfrac{\pi(a^\ast)\mu}{3r}\,(1+o(1))$; we state this only
as a first-order approximation, the exact figure being the finite sum
above.

\emph{Feedback.} A steady bias $\mu$ scales every floor on the pair by
$1-\mu$, so any guarantee whose benchmark input is the public-floor spread
$G=\max_{i,S:\,C_S>0}V_i(S)/C_S$ becomes $G/(1-\mu)$; the coverage law transforms
jointly,
$(G,\alpha)\mapsto(G/(1-\mu),\min\{1,\alpha/(1-\mu)\})$. Structural
production complementarity, directed-pair complexity, and produced-welfare
PoA are unaffected.\hfill$\square$
\end{proof}

\section{Deployment Diagnostics}\label{app:empirical}

\begin{center}
\footnotesize
\begin{minipage}{0.92\textwidth}
\centering
\textbf{Table 2.} Deployment diagnostics used only as motivation. The
visible spread is auction-level and observable; it is not the structural
complementarity factor $g$.
\medskip

\begin{tabular}{@{}ll@{}}
\toprule
Diagnostic & Value or interpretation\\
\midrule
Competitions in trace & 802,816 post-adoption solver competitions\\
Covered notional / gap & \$54.69B / \$5.13M (about 0.94 bps)\\
Visible spread tail & $\widehat G^{\mathrm{obs}}>2$: 25.3\%;
  $\widehat G^{\mathrm{obs}}>4$: 20.2\%\\
Small-notional tail & $\widehat G^{\mathrm{obs}}>2$: 63.1\%\\
Directed-pair support $P$ & directly observable from requested/paid pairs\\
Naive local reserve & 97.1\% vulnerable cell-windows in diagnostic\\
\bottomrule
\end{tabular}
\end{minipage}
\end{center}

\paragraph{Sample and filters.}
The trace contains 802,816 post-adoption solver competitions over an
eight-month window. The parser retains protocol-valid, non-filtered,
positive-score solver submissions. Full-sample quantities such as visible
spread and effective solver count use all parsed auctions. USD-denominated
quantities use smaller denominators because score conversion, notional
coverage, and sanity filters are available only for subsets of the trace.
These denominator changes are coverage restrictions, not inconsistencies.

\paragraph{Score variables.}
For each competition, $\mathrm{Score}_1$ is the winning solution score.
The reference score $\mathrm{Score}^{\mathrm{other}}_2$ is the best
positive score submitted by a solver different from the winner. The parser
groups by solver identity, removes all candidate solutions submitted by
the winning solver, and then takes the maximum remaining positive valid
score; if none exists, the reference score is zero. This best-non-winning
construction is solver-level: a winning solver's second candidate solution
is not a credible outside option.

\paragraph{Visible spread and accounting gap.}
When the denominator is positive,
$\widehat G^{\mathrm{obs}}=\mathrm{Score}_1/
\mathrm{Score}^{\mathrm{other}}_2$, with infinite spread when the
reference score is zero. The reported tail shares are computed on the full
802,816-auction sample. The dollar gap converts
$\mathrm{Score}_1-\mathrm{Score}^{\mathrm{other}}_2$ using ETH/USD prices,
winsorizes the dollar gap at the first and ninety-ninth percentiles, and
sums over the sanity-filtered denominator above. The \$54.69B covered
notional is the sell-token notional over the same denominator.

\paragraph{Reserve-vulnerability diagnostic.}
The benchmark diagnostic uses cell-windows from the reserve-construction
pipeline. With baseline quantile $\tau=0.25$, a cell-window is marked
vulnerable if the maximum eligible-sample share of any single solver is at
least $\tau$, meaning that a single solver could potentially influence the
local reserve quantile if its samples were contaminated. In the diagnostic,
97.1\% of auction-weighted cell-windows are vulnerable. The design
implication is to inherit from parent cells or use audit-only/soft reserves
when local support is sparse.
